% !TeX root = ../../金融经济学_王江_讲义及习题.tex
% 知识点：references/金融经济学/金融经济学-王江_261013.pdf，印刷页 213–223（PDF 页 228–238）
% 习题：references/金融经济学/史上最全最完整的金融经济学答案(王江版).pdf，页 76–79（印刷页 75–78）

\lecture{套利定价理论}

\ifdefined\WJHandoutKnowledge

CAPM 用市场组合的贝塔衡量系统风险。本章先不指定风险因子的经济含义，而假设收益可以由少数共同因子与剩余风险分解；再研究无套利给预期收益施加的限制。\textbf{精确复制}与\textbf{大市场中的近似复制}会得到不同强度的结论。

\begin{tcolorbox}[knowledge,title=14.1 资产收益的因子结构]
设 $N$ 只资产有 $F$ 个共同因子，收益满足线性因子模型
\[
 \widetilde r_n=\bar r_n+\sum_{k=1}^{F}b_{n,k}\widetilde f_k
 +\widetilde\varepsilon_n,\qquad n=1,\ldots,N.
 \tag{14.1}
\]
$\mathrm E[\widetilde f_k]=\mathrm E[\widetilde\varepsilon_n]=0$；$b_{n,k}$ 是资产 $n$ 对因子 $k$ 的\textbf{载荷}，$\widetilde\varepsilon_n$ 是因子无法解释的剩余风险。把因子标准化为 $\operatorname{Cov}(\widetilde f)=I$，并令因子与剩余风险不相关，则向量式为
\[
 \widetilde r=\bar r+B\widetilde f+\widetilde\varepsilon,
 \qquad \operatorname{Cov}(\widetilde r)=BB^{\mathsf T}
 +\operatorname{Cov}(\widetilde\varepsilon).
 \tag{14.2}
\]
共同因子解释跨资产共同波动；剩余风险能否充分分散，取决于它在不同资产间的协方差结构。因子可经可逆线性变换重新表示，所以单个因子的名称和载荷并不天然唯一；定价关系在相应变换后保持一致。
\end{tcolorbox}

\begin{tcolorbox}[knowledge,title=14.2 精确因子模型与无套利]
先假设无剩余风险，只有一个因子：$\widetilde r_n=\bar r_n+b_n\widetilde f$。若有两只资产的载荷不同，适当的多空组合可以使因子载荷为 0，从而得到无风险组合；也可构造载荷为 1 的\textbf{因子组合}。若无套利，所有无风险组合的收益须相同，且拥有同样因子载荷的组合须同价。于是\textbf{定理 14.1（单因子 APT）}：
\[
 \bar r_n=r_F+b_n\lambda,\qquad\forall n,
 \tag{14.5}
\]
其中 $\lambda$ 是单位因子载荷的风险溢价。若某资产不满足上式，买入便宜的一边、卖出由无风险组合和因子组合复制的另一边，可形成零成本、无随机支付差且严格正收益的套利。

精确 $F$ 因子模型写成 $\widetilde r=\bar r+B\widetilde f$。对任何零成本且零因子载荷的组合 $z$，无套利要求 $z^{\mathsf T}\bar r=0$。线性代数因此给出\textbf{定理 14.2}：
\[
 \bar r=\lambda_0\mathbf1+B\lambda,
 \quad\text{即 }\bar r_n=\lambda_0+\sum_{k=1}^{F}b_{n,k}\lambda_k.
 \tag{14.6}
\]
若存在收益线性无关的无风险组合，$\lambda_0=r_F$；因子组合可交易时，$\lambda_k$ 是第 $k$ 个标准化因子组合超出无风险利率的期望收益。因子协方差矩阵为单位阵只是标准化选择，不是经济上必须相互独立。
\end{tcolorbox}

\begin{tcolorbox}[knowledge,title=14.3 剩余风险、极限套利与近似 APT]
若每只资产还带有 $\widetilde\varepsilon_n$，有限资产的因子组合通常\textbf{不能精确复制}单只资产：消去共同因子后仍留下特殊风险。因此 (14.6) 不再由有限市场的简单复制严格推出。

考察越来越大的市场。\textbf{定义 14.1}\quad 若存在零成本组合序列，使其预期支付持续增大，而支付方差趋于 0，就称存在\textbf{极限套利}。这一概念强于只排除每个有限市场中的常规无套利机会。

\textbf{定理 14.3}\quad 若收益由固定数目 $F$ 的共同因子与有界、可充分分散的剩余风险描述，且不存在极限套利，则预期收益向量相对于某个 $F$ 因子线性定价向量的\textbf{平方偏差总和有界}：
\[
 \sum_{i=1}^{n}
 (\bar r_i^{(n)}-\lambda_0^{(n)}-b_i^{(n)\mathsf T}\lambda^{(n)})^2
 \le A,
\]
其中 $A$ 与资产数 $n$ 无关。证明的要点是把预期收益投影到由常数与因子载荷张成的空间；若正交残差的平方范数无界，可放大一个零成本、零因子载荷的组合，使其期望收益趋于无穷而剩余风险方差趋于零，违反无极限套利。

\textbf{定理 14.4}\quad 上述平方和界意味着：对任意固定 $\epsilon>0$，定价误差超过 $\epsilon$ 的资产数量至多为 $A/\epsilon^2$。故随着资产数增加，定价式对\textbf{大多数资产}近似成立；不能说每一只资产都被精确定价。

\textbf{边界例子。}即使每个有限市场没有常规套利，若越来越多独立特殊风险资产的预期超额收益保持正值，组合可分散其风险却保留收益，从而形成极限套利。无极限套利不是普通无套利的自动推论，也不是市场均衡的完整条件。
\end{tcolorbox}

\begin{tcolorbox}[knowledge,title=14.5 APT 作为均衡结果]
前述 APT 主要由支付结构和无套利出发。要说明该关系是\textbf{均衡}价格，还需偏好、禀赋、可交易性和出清等条件。\textbf{定理 14.5}给出一组充分条件：证券支付具有因子结构；无风险支付、各共同因子及参与者的禀赋可由证券支付线性表出；参与者严格风险厌恶；均衡配置为内点。在这些条件下，均衡中的个体消费可在因子支付空间内表示，个体边际效用定价条件给每个因子组合赋予溢价 $\lambda_k$，从而推出 (14.6)。

这一结果说明精确 APT 可与均衡相容，但其条件比“因子模型加无套利”更强。与 CAPM 相比，APT 容许多个系统因子，却不直接告诉我们因子的经济身份或风险溢价数值；这些仍需模型设定或数据识别。
\end{tcolorbox}

\begin{tcolorbox}[summary,title=本章小结]
精确因子模型中，因子支付可复制，常规无套利给出精确的线性预期收益关系。含剩余风险的大市场中，无极限套利只保证多数资产的近似定价，可能存在少数偏差较大的资产。若还要把 APT 解释成均衡结果，必须补充偏好、可交易支付与市场出清的条件。
\end{tcolorbox}

\fi

% !TeX root = ../../金融经济学_王江_讲义及习题.tex
% 14_APT 课后习题（由章节文件拆分）

\begin{problem}
\problemno{14.1} 市场上有无穷多的证券。它们的收益率由单因子模型给定：
\[
    \tilde{r}_{n}=\bar{r}_{n}+b_{n}\tilde{f}+\tilde{u}_{n},\quad n=1,2,\ldots
\]
这里 $\mathrm{E}[\tilde{f}]=0$；$\forall\,n,n'$，$\mathrm{E}[\tilde{u}_{n}]=0$，$\mathrm{E}[\tilde{f}\tilde{u}_{n}]=0$；$\operatorname{Var}[\tilde{u}_{n}]\le\sigma^{2}$；以及 $\mathrm{E}[\tilde{u}_{n}\tilde{u}_{n'}]=0$。一个组合由 $\{w_{1},w_{2},\ldots\}$ 给定。如当 $N$ 很大时，
\[
    \sum_{n}^{N}w_{n}^{2}\simeq O(N^{-1})
\]
则称它是充分分散的，其中 $N$ 代表组合中绝对权重最大的股票数目而 $O(\cdot)$ 表示同阶或更高阶。
\begin{enumerate}[label=(\alph*)]
    \item 计算充分分散组合的方差。
    \item 证明它不受单个证券的特殊风险的影响。
    \item 更一般地讨论什么因素决定了充分分散组合的风险，这里风险的度量是收益率的方差。
\end{enumerate}
\end{problem}
\begin{solution}
\begin{enumerate}[label=(\alph*)]
    \item 组合收益率为
    \[
        \tilde{r}_{p}=\sum w_{n}\tilde{r}_{n}=\sum w_{n}\bar{r}_{n}+\Bigl(\sum w_{n}b_{n}\Bigr)\tilde{f}+\sum w_{n}\tilde{u}_{n}
    \]
    组合风险为
    \[
        \operatorname{Var}(\tilde{r}_{p})=\Bigl(\sum w_{n}b_{n}\Bigr)^{2}\sigma_{f}^{2}+\sum w_{n}^{2}\operatorname{Var}(\tilde{u}_{n})
    \]
    因为组合是充分分散的，故
    \[
        \sum w_{n}^{2}\operatorname{Var}(\tilde{u}_{n})\le\sigma^{2}\sum w_{n}^{2}\to 0,\quad N\to\infty
    \]
    因此
    \[
        \operatorname{Var}(\tilde{r}_{p})=\Bigl(\sum w_{n}b_{n}\Bigr)^{2}\sigma_{f}^{2}
    \]
    \item $\operatorname{Var}(\tilde{r}_{p})$ 中没有 $\operatorname{Var}(\tilde{u}_{n})$ 项，即组合风险不受特殊风险的影响；
    \item 影响充分分散组合的风险有两个：一是因子的风险，另一个是组合的载荷。
\end{enumerate}
\end{solution}

\begin{problem}
\problemno{14.2} 假设资产 $A$、$B$、$C$ 的收益服从精确两因子模型：
\begin{align}
    r_{A}&=a_{A}+b_{A,1}\tilde{f}_{1}+b_{A,2}\tilde{f}_{2}\tag{14.1}\\
    r_{B}&=a_{B}+b_{B,1}\tilde{f}_{1}+b_{B,2}\tilde{f}_{2}\tag{14.2}\\
    r_{C}&=a_{C}+b_{C,1}\tilde{f}_{1}+b_{C,2}\tilde{f}_{2}\tag{14.3}
\end{align}
其中 $\tilde{f}_{1}$ 和 $\tilde{f}_{2}$ 是两个均值为 0、方差为 1 的因子。$r_{F}$ 是无风险利率。为消除套利机会，$a$、$b$ 以及 $r_{F}$ 应该满足什么条件？
\end{problem}
\begin{solution}
构造一个组合，组合权重为 $\begin{bmatrix}x_{1} \\ x_{2} \\ 1-x_{1}-x_{2}\end{bmatrix}$，那么组合的收益率为
\[
    \begin{aligned}
        \tilde{r}_{p}&=\bigl[(a_{A}-a_{C})x_{1}+(a_{B}-a_{C})x_{2}\bigr]+\bigl[(b_{A,1}-b_{C,1})x_{1}+(b_{B,1}-b_{C,1})x_{2}\bigr]\tilde{f}_{1}\\
        &\quad+\bigl[(b_{A,2}-b_{C,2})x_{1}+(b_{B,2}-b_{C,2})x_{2}\bigr]\tilde{f}_{2}
    \end{aligned}
\]
如果不存在套利机会，那么
\[
    \begin{aligned}
        (b_{A,1}-b_{C,1})x_{1}+(b_{B,1}-b_{C,1})x_{2}&=0\\
        (b_{A,2}-b_{C,2})x_{1}+(b_{B,2}-b_{C,2})x_{2}&=0\\
        (a_{A}-a_{C})x_{1}+(a_{B}-a_{C})x_{2}&=r_{F}
    \end{aligned}
\]
也就是说，若 $x\in\{y\in\mathrm{R}^{2}:y\perp\begin{bmatrix}b_{A,1}-b_{C,1} \\ b_{B,1}-b_{C,1}\end{bmatrix},\,y\perp\begin{bmatrix}b_{A,2}-b_{C,2} \\ b_{B,2}-b_{C,2}\end{bmatrix}\}$，那么 $x^{\mathrm{T}}\begin{bmatrix}a_{A}-a_{C} \\ a_{B}-a_{C}\end{bmatrix}=r_{F}$。
\end{solution}

\begin{problem}
\problemno{14.3} 考虑一个有无穷多资产和 $K$ 个因子的因子模型：
\[
    \tilde{r}_{n}=\bar{r}_{n}+b_{n}\tilde{f}+\tilde{u}_{n},\quad n=1,2,\ldots
\]
其中 $\tilde{f}=\begin{bmatrix}\tilde{f}_{1} \\ \ldots \\ \tilde{f}_{K}\end{bmatrix}$ 是因子向量，$b_{n}=\begin{bmatrix}b_{n,1} \\ \ldots \\ b_{n,K}\end{bmatrix}$ 资产 $n$ 的因子载荷向量，而 $\forall\,n,n'$，$\mathrm{E}[\tilde{f}]=0$，$\mathrm{E}[\tilde{u}_{n}]=0$，$\mathrm{E}[\tilde{f}\tilde{f}^{\mathrm{T}}]=I$，$\mathrm{E}[\tilde{f}\tilde{u}_{n}]=0$，$\mathrm{E}[\tilde{u}_{n}\tilde{u}_{n'}]=0$，其中 $I$ 是单位阵。考虑另外一组因子 $\tilde{f}'$，它们是 $\tilde{f}$ 的线性变换：
\[
    \tilde{f}'=\mathrm{A}\tilde{f}
\]
这里 $\mathrm{A}$ 是一个 $(K\times K)$ 的矩阵。
\begin{enumerate}[label=(\alph*)]
    \item 证明存在 $\tilde{f}'$ 使得
    \[
        \tilde{r}_{n}=\bar{r}_{n}+b_{n}\tilde{f}'+\tilde{u}_{n},\quad n=1,2,\ldots
    \]
    其中 $\mathrm{E}[\tilde{f}']=0$，$\mathrm{E}[\tilde{f}'\tilde{f}'^{\mathrm{T}}]=I$，$\mathrm{E}[\tilde{f}'\tilde{u}_{n}]=0$，$\forall\,n$。因此，因子模型不是唯一的。
    \item 证明因子模型有许多因子不同但却等价的表述形式。
\end{enumerate}
\end{problem}
\begin{solution}
\begin{enumerate}[label=(\alph*)]
    \item 若 $A$ 为正交阵（即 $A^{\mathrm{T}}A=I$），那么
    \[
        \begin{aligned}
            \tilde{r}_{n}&=\bar{r}_{n}+b_{n}^{\mathrm{T}}\tilde{f}+\tilde{u}_{n}\\
            &=\bar{r}_{n}+b_{n}^{\mathrm{T}}A^{\mathrm{T}}A\tilde{f}+\tilde{u}_{n}\\
            &=\bar{r}_{n}+(Ab_{n})^{\mathrm{T}}(A\tilde{f})+\tilde{u}_{n}\\
            &=\bar{r}_{n}+b_{n}'^{\mathrm{T}}\tilde{f}'+\tilde{u}_{n}
        \end{aligned}
    \]
    可以验证：
    \[
        \mathrm{E}[\tilde{f}']=0,\quad\mathrm{E}[\tilde{f}'\tilde{f}'^{\mathrm{T}}]=I,\quad\mathrm{E}[\tilde{f}'\tilde{u}_{n}]=0
    \]
    \item 因为 $K$ 阶正交阵有很多，因而等价的表述形式也有很多。
\end{enumerate}
\end{solution}

\begin{problem}
\problemno{14.4*} 给定练习 14.3 描述的因子模型，APT 要求
\[
    \bar{r}_{n}-r_{F}=\lambda^{\mathrm{T}}b_{n}
\]
其中 $\lambda=\begin{bmatrix}\lambda_{1} \\ \lambda_{2} \\ \ldots \\ \lambda_{K}\end{bmatrix}$ 是因子风险溢价向量。证明存在一个等价的因子模型，它的因子向量为 $\tilde{f}^{*}$ 而对应的因子风险溢价向量为 $\lambda^{*}=\begin{bmatrix}\lambda_{1}^{*} \\ 0 \\ \ldots \\ 0\end{bmatrix}$，即只有第一个因子有风险溢价而其它因子虽然代表了不可分散风险，却没有风险溢价。
\end{problem}
\begin{solution}
考虑这样的正交阵：$A=\bigl[\frac{1}{\sqrt{\lambda^{\mathrm{T}}\lambda}}\lambda,\,U\bigr]$，其中
\[
    U=[U_{1},\,U_{2},\,\cdots,\,U_{K-1}],\quad U_{i}^{\mathrm{T}}U_{j}=\delta_{ij},\quad U_{i}^{\mathrm{T}}\lambda=0
\]
\[
    \delta_{ij}=\begin{cases}
        1&i=j\\
        0&i\neq j
    \end{cases}
\]
可以验证：$AA^{\mathrm{T}}=I$，即 $A$ 是正交阵，且 $A^{\mathrm{T}}\lambda=\begin{bmatrix}\sqrt{\lambda^{\mathrm{T}}\lambda} \\ \,0 \\ \,0 \\ \,\cdots \\ \,0\end{bmatrix}$，那么
\[
    \begin{aligned}
        \tilde{r}_{n}&=\bar{r}_{n}+b_{n}^{\mathrm{T}}AA^{\mathrm{T}}\tilde{f}+\tilde{u}_{n}\\
        &=\bar{r}_{n}+(A^{\mathrm{T}}b_{n})^{\mathrm{T}}(A^{\mathrm{T}}\tilde{f})+\tilde{u}_{n}\\
        &=\bar{r}_{n}+b_{n}^{*\,\mathrm{T}}\tilde{f}^{*}+\tilde{u}_{n}
    \end{aligned}
\]
是一个等价表述。并且 $\bar{r}_{n}-r_{F}=\lambda^{*\,\mathrm{T}}b_{n}^{*}$，其中，$\lambda^{*}=A^{\mathrm{T}}\lambda=\begin{bmatrix}\sqrt{\lambda^{\mathrm{T}}\lambda} \\ \,0 \\ \,0 \\ \,\cdots \\ \,0\end{bmatrix}$。
\end{solution}


